% ============================================================
%  Research Design Chapter Content
%  Research Position, Methodology, Methods, Ethics
%  Henry Kabuye PhD Thesis — Northumbria University
% ============================================================

\lettrine{E}{stablishing} a coherent research design is foundational to any doctoral investigation.
The methodological and philosophical choices made at the outset of a research programme determine
the types of knowledge it produces, the criteria by which its findings are evaluated, and the degree
to which its conclusions can be generalised beyond the specific experimental context in which they
were obtained \cite{saunders2019research}.
This chapter articulates the philosophical stance and methodological framework underpinning the
three empirical investigations presented in Chapters~3, 4, and 5.
Each investigation builds incrementally upon its predecessor: the first establishes a single-agent
uncertainty-aware detection baseline \cite{kabuye2025explainable}; the second extends this to a
two-agent multi-modal framework with uncertainty-driven triage \cite{kabuye2025sabre_eswa}; and
the third introduces a knowledge-grounded third agent with inter-agent consultation and continual
learning.
All three studies share the same philosophical orientation, experimental strategy, and evaluation
framework, which are described here to avoid repetition and to demonstrate the internal consistency
of the research programme as a whole.

% ============================================================
\section{Research Position}
\label{sec:research_position}
% ============================================================

The philosophical orientation of a research programme is most systematically articulated through
the Research Onion framework proposed by Saunders, Lewis and Thornhill \cite{saunders2019research}.
Working from the outermost philosophical layer inward through research approach, strategy, time
horizon, and data collection and analysis techniques, the framework ensures that each methodological
choice is coherent with the philosophical commitments from which it derives.
Figure~\ref{fig:research-onion} illustrates the Research Onion and the position adopted within each
layer for this thesis.
The three innermost layers — strategy, time horizon, and techniques — are addressed in
Section~\ref{sec:methodology} and Section~\ref{sec:methods}; the three outermost philosophical
layers are addressed below.

\begin{figure}[ht]
    \centering
    \includegraphics[width=0.78\textwidth]{ResearchOnion}
    \caption{The Research Onion (adapted from \protect\cite{saunders2019research}), illustrating the
    nested layers of philosophical and methodological choice.
    This thesis occupies a post-positivist realist position, adopts a deductive approach,
    employs an experimental mono-method quantitative strategy, and applies a cross-sectional
    time horizon with a simulated longitudinal component in Chapter~5.}
    \label{fig:research-onion}
\end{figure}

% ------------------------------------------------------------
\subsection{Ontology}
\label{sec:ontology}
% ------------------------------------------------------------

Ontology concerns the nature of reality and what entities exist in the world independently of an
observer's perceptions of them.
The ontological position adopted in this thesis is \emph{post-positivist realism}: the view that
ransomware constitutes an objective, observer-independent phenomenon with measurable behavioural and
structural signatures.
A ransomware binary either encrypts user data or it does not; it either contacts
command-and-control infrastructure or it does not; it either exhibits the MITRE ATT\&CK tactic
patterns documented by practitioner communities or it does not.
These are facts about the world, not social constructions, and they are in principle discoverable
through careful measurement \cite{saunders2019research}.

This realist stance does not, however, imply that all knowledge about ransomware is complete or
certain.
The thesis departs from na\"{i}ve positivism in acknowledging that measurements are inherently noisy,
datasets are finite, and the classification boundary between malicious and benign software is
probabilistic in marginal cases.
Novel ransomware families can exhibit structural and behavioural signatures that overlap
substantially with legitimate software, particularly during the reconnaissance and staging phases
of an attack that precede the encryption payload \cite{razaulla2023age, beaman2021ransomware}.
The precise boundary of the malicious class is therefore epistemically uncertain even if its
existence is ontologically real.

This observation directly motivates the epistemic uncertainty quantification that is central to
each empirical chapter in this thesis.
Rather than merely acknowledging measurement uncertainty as a limitation to be minimised, the
research operationalises it as an executive control variable embedded directly in the detection
pipeline's triage logic.
The position is consistent with what Saunders et al.\ \cite{saunders2019research} term
\emph{critical realism}: an ontological stance that accepts the independent existence of a real
world but holds that knowledge of it is always mediated by observation, theory, and the instruments
of measurement.
In applied security research, this position has practical consequences: a model that assigns
overconfident posterior probabilities to borderline cases provides less actionable intelligence to
an analyst than one that correctly identifies and communicates the limits of its own
knowledge \cite{gawlikowski2023Uncertainity, Doan2024uncertainBay}.

% ------------------------------------------------------------
\subsection{Epistemology}
\label{sec:epistemology}
% ------------------------------------------------------------

Epistemology concerns how knowledge is acquired, validated, and distinguished from belief or
conjecture.
The epistemological stance adopted here is broadly \emph{empiricist} and quantitative: knowledge
about ransomware detection is generated by designing controlled experiments, collecting labelled
data from authoritative sources, training models on that data, and evaluating their predictions
against held-out ground truth.
The primary epistemological commitment is to \emph{falsifiability}: every empirical claim in this
thesis is, in principle, refutable by experiment.
The proposition that an uncertainty-aware triage policy reduces false containment events compared to
a threshold-only policy, for example, is directly testable by comparing the two policies' decisions
against ground-truth labels on a held-out evaluation partition.

The thesis operates within a \emph{post-positivist} paradigm rather than a strictly positivist one.
Strict positivism holds that scientific knowledge is derived from direct, value-free observation and
that measurement produces certain knowledge about the world.
Post-positivism acknowledges that all observations are theory-laden, that all measurement is
imperfect, and that probabilistic rather than deterministic claims are appropriate when working with
complex natural phenomena \cite{saunders2019research}.
This is directly reflected in the methodology: each model is evaluated not only by point-estimate
accuracy metrics but also by calibration diagnostics that assess the reliability of probability
estimates across the full confidence distribution \cite{guo2017calibration}, and by escalation-rate
analyses that quantify the operational cost of residual uncertainty.

Ground truth for this research is established through expert-consensus mechanisms.
The RanSAP and RDset benchmarks provide labels derived from multi-scanner analysis and controlled
dynamic sandbox execution; the RansKnow dataset provides labels derived from practitioner-validated
thresholds on ATT\&CK tactic prevalence.
These are accepted proxies for ground truth rather than perfect knowledge of it — consistent with
the post-positivist acknowledgement that all empirical knowledge is provisional and subject to
revision as better data and methods become available.

% ------------------------------------------------------------
\subsection{Axiology}
\label{sec:axiology}
% ------------------------------------------------------------

Axiology concerns the role of values in the design and conduct of research, including both the
values of the investigator and the values embedded in the choice of research problem.
The axiological position adopted in this thesis is broadly \emph{objectivist}: the investigator
aims to maintain value-neutrality in data collection, model training, and the reporting of
evaluation results, disclosing performance honestly regardless of whether findings support or
contradict prior hypotheses.
All evaluation metrics — including cases where the proposed frameworks fall short of stated targets
or produce results that require qualification — are reported in full in the empirical chapters.

However, the choice of research problem is itself a value-laden act.
This thesis is motivated by the measurable harm that ransomware causes to public-sector
organisations, critical infrastructure, and individual users \cite{europol2020iocta,
beaman2021ransomware}.
The research is therefore oriented toward the development of defensive detection capabilities rather
than toward understanding ransomware for its own sake or producing knowledge that could assist
adversarial actors.
This orientation does not compromise methodological objectivity but it does delimit the scope of
the investigation: techniques that might yield detection insights but require active deployment of
malicious payloads in non-sandboxed environments, reverse engineering of proprietary production
systems, or collection of personal data from affected individuals fall outside the scope of this
research.

All datasets used are either publicly available research benchmarks or researcher-constructed
from open-source public materials.
No personally identifiable information was collected or processed at any stage of the investigation.
These commitments are addressed in greater detail in Section~\ref{sec:ethics}.

% ============================================================
\section{Research Methodology}
\label{sec:methodology}
% ============================================================

Having established the philosophical position of the research, this section addresses the three
inner layers of the Research Onion: the research approach, the research strategy, and the time
horizon.
Together these choices define the logical structure of each empirical investigation and ensure
that the methods employed in Chapters~3--5 are coherent with the post-positivist realist stance
described above.

% ------------------------------------------------------------
\subsection{Research Approach}
\label{sec:approach}
% ------------------------------------------------------------

The research approach adopted throughout this thesis is \emph{deductive}.
A deductive approach begins with general theory and established prior knowledge, derives specific
testable hypotheses from that theory, and then subjects those hypotheses to empirical
evaluation \cite{saunders2019research}.
The general theoretical premises of this thesis are drawn from the established literature on
machine learning-based malware detection \cite{razaulla2023age, gibert2020rise}, uncertainty
quantification in deep neural networks \cite{gal2016dropout, gawlikowski2023Uncertainity},
neuro-symbolic artificial intelligence \cite{Besold2021neurosymbolicsurvey}, and multi-agent
detection systems \cite{MultiAgentAdaptiveDetection2023}.
From these premises, the thesis derives and tests a progression of three hypotheses:

\begin{enumerate}
    \item That integrating data augmentation, Monte Carlo Dropout uncertainty quantification, and
    post-hoc explainability within a single-agent CNN framework yields detection accuracy and
    analyst interpretability superior to baseline classifiers on the RanSAP and RDset benchmarks.

    \item That decomposing detection across modality-specialised agents fused through a learned
    decision layer, with an uncertainty-driven triage policy governing the balance between
    autonomous containment and human escalation, yields superior operational control compared to
    a single unified model under both safety-optimal and low-escalation operating conditions.

    \item That augmenting the two-agent framework with a knowledge-grounded third agent trained
    on threat-intelligence data, a Peer Consultation Protocol enabling inter-agent uncertainty
    communication, and an Expert Feedback Loop for continual learning yields robust detection
    performance under distributional shift introduced by novel ransomware family variants.
\end{enumerate}

Each hypothesis is tested in a separate empirical chapter using controlled experiments on the
corresponding datasets, with performance assessed against the evaluation framework described in
Section~\ref{sec:eval_framework}.

% ------------------------------------------------------------
\subsection{Research Strategy}
\label{sec:strategy}
% ------------------------------------------------------------

The research strategy employed is \emph{experimental}.
Controlled experiments on benchmark datasets with pre-defined and fixed training, validation, and
test partitions allow direct, reproducible comparison between the proposed frameworks and
established baseline approaches \cite{saunders2019research}.
Each empirical chapter follows a consistent experimental protocol: the full dataset is partitioned
before any model training; hyperparameter selection is performed exclusively on the validation
partition through Bayesian optimisation \cite{optuna2019}; and all final performance figures are
computed on the held-out test partition only, which is never used for model selection or threshold
tuning.

The experimental strategy supports two complementary forms of analysis.
\emph{Confirmatory analysis} tests the specific performance hypotheses stated above using
pre-specified evaluation metrics.
\emph{Exploratory analysis} examines uncertainty distributions, calibration curves,
counterfactual explanations, and the effects of the peer consultation mechanism in order to
characterise model behaviour and generate insights that motivate the extensions presented in
subsequent chapters.
Ablation studies are employed in each chapter to isolate the contribution of individual
framework components — data augmentation, MC Dropout, the triage policy thresholds, peer
consultation, and continual learning — from the performance of the system as a whole.

A \emph{mono-method quantitative} methodological choice is adopted throughout.
All data are numeric, all hypotheses are stated in quantitative terms, and all findings are
expressed as measurable performance metrics.
Mixed-methods or qualitative approaches are not employed, as the research questions are
amenable to precise quantitative formulation and the available datasets provide sufficient
quantitative evidence to answer them.

% ------------------------------------------------------------
\subsection{Time Horizon}
\label{sec:time_horizon}
% ------------------------------------------------------------

The primary time horizon of the research is \emph{cross-sectional}: all three empirical
investigations evaluate their respective frameworks on static, fixed datasets at a single point
in time, reflecting the standard practice in machine learning benchmarking
research \cite{saunders2019research}.
Cross-sectional evaluation on held-out test sets provides the most direct and reproducible
measure of a model's generalisation ability, and is the standard adopted by the security
machine learning community.

The third empirical study (Chapter~5) introduces a controlled \emph{longitudinal element} through
the Expert Feedback Loop mechanism, which models the behaviour of a detection framework receiving
labelled analyst feedback over successive deployment cycles and updating its decision logic
incrementally through experience replay \cite{rolnick2019experience}.
This element is evaluated through a simulated sequential replay protocol on the RansKnow
dataset rather than through live deployment, preserving the experimental control of a
cross-sectional design while demonstrating the system's capacity for temporal adaptation under
concept drift.
The distinction between this simulated longitudinal evaluation and genuine longitudinal study
in a production environment is discussed in Chapter~5.

% ============================================================
\section{Research Methods and Techniques}
\label{sec:methods}
% ============================================================

The preceding sections established the philosophical and strategic orientation of the research.
This section describes the specific datasets, model architectures, and evaluation procedures
through which the three research hypotheses are pursued.

% ------------------------------------------------------------
\subsection{Datasets and Data Sources}
\label{sec:datasets}
% ------------------------------------------------------------

Three datasets provide the empirical basis for the investigations in this thesis.
Each captures a distinct dimension of ransomware activity, reflecting the thesis's core
argument that ransomware exhibits multiple observable signatures most effectively detected by
specialised models trained on different data modalities.
The key characteristics of all three datasets are summarised in Table~\ref{tab:datasets}.

\textbf{RanSAP} (Ransomware System API behaviour dataset) \cite{hirano2022ransap} captures
dynamic runtime behaviour through system API call sequences recorded during controlled sandbox
execution of ransomware samples and benign applications.
The dataset records 30-second time-window API traces across multiple ransomware families,
providing a temporal signal that reflects the characteristic I/O access patterns, file
enumeration sequences, and encryption activity exhibited during a live ransomware attack.
In this thesis, RanSAP serves as the primary dataset for training the Behavioural Agent in
Chapters~4 and~5, and also contributes to the multi-class classification task in Chapter~3.

\textbf{RDset} provides PE-format static analysis features derived from portable executable
file headers, including section entropy profiles, import address table contents, and PE
flag configurations.
These features encode the semantic structure of executable files without requiring dynamic
execution, enabling rapid pre-execution screening of samples at the point of entry.
The Semantic Agent in Chapters~4 and~5 is trained on the RDset, and Chapter~3 employs these
features in the CNN-based binary classifier.

\textbf{RansKnow} is a researcher-constructed threat-intelligence dataset comprising 307
samples sourced from publicly available YouTube security briefings, conference presentations,
and practitioner walkthroughs.
Each sample is represented by 23 features spanning ten MITRE ATT\&CK tactic occurrence
indicators, eight ransomware-associated tool count features, three platform-specificity
signals, and two aggregate totals.
Binary labels are derived from the operational threshold condition
(\emph{Tactic\textunderscore Total}~$>$~10) OR (\emph{Tool\textunderscore Total}~$>$~2),
which reflects practitioner consensus on the minimum tactical breadth and tooling complexity
that distinguishes a documented ransomware deployment from benign or marginal activity.
RansKnow is used in Chapter~5 to train the Knowledge Agent.

\begin{table}[ht]
\centering
\caption{Summary characteristics of the three datasets used in this thesis.}
\label{tab:datasets}
\begin{tabular}{@{}llll@{}}
\toprule
\textbf{Dataset}  & \textbf{Modality}         & \textbf{Primary Use}              & \textbf{Chapter} \\
\midrule
RanSAP   & Runtime API traces        & Behavioural Agent training        & 3, 4, 5          \\
RDset    & PE / static file features & Semantic Agent training           & 3, 4, 5          \\
RansKnow & Threat-intelligence       & Knowledge Agent training          & 5                \\
\bottomrule
\end{tabular}
\end{table}

% ------------------------------------------------------------
\subsection{Machine Learning Architecture and Training}
\label{sec:ml_arch}
% ------------------------------------------------------------

The primary classification architecture employed across all empirical chapters is a
\emph{Convolutional Neural Network} (CNN) adapted for structured tabular feature representations.
CNNs have been demonstrated to extract effective local feature interactions from security
data \cite{kolosnjaji2016deep, zahoora2022Deep}, and their convolutional structure provides
a natural basis for applying Monte Carlo Dropout uncertainty estimation without modification
to the standard training procedure \cite{gal2016dropout}.
Each CNN backbone consists of stacked convolutional blocks followed by fully connected layers;
specific architectural hyperparameters are reported within each empirical chapter.

\emph{Uncertainty quantification} is implemented via Monte Carlo (MC) Dropout throughout all
three studies.
At inference time, dropout masks are sampled $T$ times per input, producing a distribution over
class probabilities from which the mean prediction $\bar{p}$ and predictive standard deviation
$\sigma$ are computed.
This approach yields calibrated epistemic uncertainty estimates that reflect genuine uncertainty
about the decision boundary, particularly on inputs near or between the training
classes \cite{gawlikowski2023Uncertainity}.
The value of $T$ is tuned per study to balance uncertainty quality against inference latency,
with values of $T=30$ used in Chapter~3 and Chapter~4, and an adaptive mechanism
($T_{\min}=5$, $T_{\max}=50$) employed by the Knowledge Agent in Chapter~5.

\emph{Score-space data augmentation} via Conditional Tabular GAN (CTGAN) \cite{xu2019ctgan} is
applied in Chapters~4 and~5 to augment the training distribution of the Decision MLP on the
fused agent score vectors.
CTGAN learns the conditional distribution of features given class labels through an adversarial
training procedure and generates high-fidelity synthetic tabular samples that preserve feature
correlations without overfitting to the artefacts of a finite dataset.
This is a departure from the raw-feature augmentation employed in Chapter~3 and is justified
by the substantially lower dimensionality of the score space (four dimensions in Chapter~4,
six dimensions in Chapter~5) compared to the original feature representations.

\emph{Explainability} is provided through SHAP (SHapley Additive exPlanations) feature
attribution values \cite{lundberg2017shap} and, in Chapter~4, through a novel counterfactual
explanation protocol that identifies the minimal feature perturbation necessary to reverse a
ransomware prediction.
These mechanisms address the practical analyst requirement for interpretable explanations that
accompany automated detection decisions and can be acted upon without specialist knowledge of
the underlying model architecture.

\emph{Temperature calibration} is applied to all agent posterior probabilities following
\cite{guo2017calibration} before fusion, ensuring that each agent's confidence score reflects
a reliable frequency-of-correctness estimate rather than an uncalibrated logit output.
The calibration temperature is fitted by minimising the negative log-likelihood on a held-out
validation set and is fixed thereafter.

\emph{Continual learning} in Chapter~5 employs the Expert Feedback Loop, which updates the
Decision MLP weights after each labelled feedback batch through reservoir-sampled experience
replay \cite{rolnick2019experience}, preventing catastrophic forgetting of previously seen
distributions while incorporating new information from analyst corrections.
Only the Decision MLP (approximately 4,000 parameters) is updated; all agent backbone CNNs
remain frozen throughout the feedback process.

% ------------------------------------------------------------
\subsection{Evaluation Framework}
\label{sec:eval_framework}
% ------------------------------------------------------------

Model performance is quantified using a complementary set of metrics chosen to characterise
different aspects of operational utility.
\emph{AUC} (Area Under the ROC Curve) provides a threshold-independent measure of
discriminative ability.
\emph{Balanced accuracy} (the arithmetic mean of per-class recall) accounts for class
imbalance without requiring threshold selection and is the primary metric used for
Bayesian hyperparameter optimisation.
\emph{Precision}, \emph{recall}, and \emph{F1-score} characterise the precision--recall
trade-off at the selected operating threshold.

\emph{Calibration quality} is assessed via the Expected Calibration Error (ECE) and
reliability diagrams following the temperature scaling protocol of Guo et al.\
\cite{guo2017calibration}, which decompose the gap between model confidence and empirical
accuracy across equal-width confidence bins.

The \emph{operational performance} of the uncertainty-aware triage policy in Chapters~4 and~5
is assessed through two complementary metrics: the \emph{escalation rate} (the fraction of
test-set predictions forwarded to human review) and the \emph{containment precision} (the
fraction of auto-contained predictions that correspond to genuine ransomware).
Together these metrics characterise the practical utility of the framework for a human analyst
who must act on the system's decisions under time pressure.

Hyperparameter selection in all three studies employs Bayesian optimisation with
Tree-structured Parzen Estimator (TPE) sampling as implemented in the Optuna
library \cite{optuna2019}, which builds probabilistic surrogate models of the objective
function and efficiently directs the search toward regions of the hyperparameter space likely
to yield high-performing configurations.

% ============================================================
\section{Ethical Considerations}
\label{sec:ethics}
% ============================================================

The research presented in this thesis raises a number of ethical considerations that were
addressed proactively at the design stage.

\textbf{Data ethics and privacy.}
All datasets used in this thesis are either publicly available research benchmarks or
researcher-constructed from openly accessible public materials.
The RanSAP and RDset benchmarks are distributed by their originators for academic use without
restriction on non-commercial research.
The RansKnow dataset was constructed exclusively from publicly available YouTube videos,
conference proceedings, practitioner blog posts, and security briefing transcripts; no private
communications, organisational records, or personally identifiable information were collected
or processed at any stage.
The research does not interact with any production systems, operational networks, or
live user data.

\textbf{Malware handling.}
The thesis does not require direct interaction with live malware payloads or command-and-control
infrastructure.
The RanSAP dataset provides pre-collected API execution traces from controlled sandbox
environments operated by the dataset's originators; this thesis uses these traces as static
feature matrices without re-execution of any malicious code.
The RDset provides pre-extracted PE structural features; no executable binaries are loaded or
analysed dynamically within this research.
All experiments were conducted on isolated research computing infrastructure with no external
network connectivity during model training or evaluation.

\textbf{No human subjects.}
The research involves no human participants, surveys, interviews, user studies, or any
collection of personal data.
No ethical approval for human-subjects research was required or sought.
The research does not profile any individual and does not produce outputs that could be used
to identify, monitor, or discriminate against any person.

\textbf{Defensive orientation and responsible disclosure.}
The multi-agent framework described in this thesis is a prototype research system evaluated
on benchmark datasets.
It is designed exclusively to detect ransomware activity and produce analyst-interpretable
explanations, and it does not produce techniques, tools, or knowledge that would provide
operational advantages to adversarial actors beyond what is already available in the public
security research literature.
Any future deployment of these methods in an operational security environment would require
independent security auditing, institutional oversight, and appropriate governance to ensure
that automated containment decisions are exercised responsibly, with suitable human review
mechanisms in place for escalated or uncertain cases.

% ============================================================
\section{Thesis Chapter Structure}
\label{sec:chapter_structure}
% ============================================================

The following provides a concise orientation to the chapters of this thesis and their
relationship to the research objectives stated above.

\textbf{Chapter~1: Introduction} provides the motivating context for the research, documents
the scale and impact of ransomware attacks on modern organisations and critical
infrastructure \cite{europol2020iocta, beaman2021ransomware}, identifies the principal
limitations of existing detection approaches with respect to uncertainty, interpretability,
and multi-modal reasoning, and states the specific research objectives and contributions of
the thesis.

\textbf{Chapter~2: Research Design} (this chapter) establishes the philosophical position,
methodological framework, dataset characteristics, evaluation approach, and ethical
commitments that underpin all three empirical investigations, ensuring that their findings
are interpreted within a coherent epistemological context.

\textbf{Chapter~3: Explainable and Uncertainty-Aware AI-based Ransomware Detection} presents
the first empirical study, published as a journal article in IEEE Access
\cite{kabuye2025explainable}.
This chapter develops a single-agent CNN framework that integrates MC Dropout uncertainty
quantification, CTGAN-based class-imbalance augmentation, and SHAP/LIME post-hoc
explainability.
The study establishes the performance baseline upon which the subsequent multi-agent
architectures build, and demonstrates that operationalising uncertainty alongside
accuracy-based metrics provides a richer characterisation of model behaviour than
accuracy metrics alone.

\textbf{Chapter~4: Agentic SABRE --- An Uncertainty-Aware Neuro-Symbolic Multi-Agent
Framework for Adaptive Ransomware Detection} presents the second empirical study
(accepted for publication in \emph{Expert Systems with Applications})
\cite{kabuye2025sabre_eswa}.
This chapter decomposes ransomware detection into two modality-specialised agents (Semantic
and Behavioural) fused through a learned Decision Agent.
The uncertainty-aware triage policy is introduced as the primary operational contribution,
enabling the system to balance autonomous containment against human escalation through
learnable threshold parameters optimised with Bayesian TPE.

\textbf{Chapter~5: Extended Agentic SABRE --- Knowledge-Augmented Three-Agent Ransomware
Detection with Peer Consultation and Continual Learning} presents the third empirical study
(currently under review at \emph{IEEE Transactions on Information Forensics and Security}).
This chapter introduces the Knowledge Agent trained on the RansKnow threat-intelligence
dataset, the Peer Consultation Protocol through which agents communicate epistemic uncertainty
before fusion, and the Expert Feedback Loop for analyst-guided continual learning.
Together these three extensions address the limitations of purely data-driven detection and
bring the framework closer to the neuro-symbolic ideal of combining learned statistical
representations with structured domain knowledge.

% ============================================================
\section{Summary}
\label{sec:resdesign_summary}
% ============================================================

This chapter has established the philosophical and methodological foundations of the research
programme.
A post-positivist realist ontology justifies the treatment of ransomware behaviour as
objectively measurable while acknowledging the inherent probabilistic uncertainty of
measurements near the decision boundary.
An empiricist epistemology grounds knowledge claims in controlled, reproducible experiments
evaluated against authoritative ground truth.
A deductive research approach derives testable hypotheses from the established literature and
tests them through mono-method quantitative experiments under a cross-sectional (with
limited longitudinal) time horizon.
Three purpose-selected datasets covering runtime API behaviour, static PE-format features,
and practitioner-curated threat intelligence provide the empirical material; CNNs with MC
Dropout uncertainty quantification, CTGAN score-space augmentation, and SHAP/LIME
explainability constitute the core technical framework; and the evaluation combines
discriminative accuracy, calibration quality, and operational triage metrics to assess each
system against both technical performance and practical deployment criteria.
The following three chapters present the empirical investigations conducted within this
framework, building from the single-agent baseline through the two-agent triage system
to the full three-agent knowledge-augmented architecture.
